\documentclass[12pt]{article}

\usepackage[T1]{fontenc}
\usepackage[utf8]{inputenc}
\usepackage{lmodern}
\usepackage[a4paper,margin=1in]{geometry}
\usepackage{setspace}
\usepackage{amsmath}
\usepackage{amssymb}
\usepackage{hyperref}
\usepackage[style=apa]{biblatex}
\DeclareLanguageMapping{english}{english-apa}

\addbibresource{refs.bib}

\title{Nevo and Rosen (2012): Simulations and Extensions}
\author{Maximilian Grotz\\Toulouse School of Economics}
\date{\today}

\onehalfspacing

\begin{document}

\maketitle

\section{Nevo and Rosen (2012)}

\section{ Nevo and Rosen (2012) as Sensitivity Analysis}

\subsection{Sensitivity Analysis: General Approach}

This material has mainly been taken from Paul Diegert's lecture slides on sensitivity analysis. 

Consider a family of models indexed by a sensitivity parameter $\delta \in [\underline{\delta}, \bar{\delta}]$:
\[
\mathcal{M}(\delta) \subseteq \mathcal{M}(\delta') \subseteq \mathcal{M}
\quad \text{for } \delta \le \delta'.
\]

The parameter $\delta$ governs the strength of assumptions: larger $\delta$ corresponds to weaker assumptions and hence a larger model class.

Assume models are \emph{complete}, i.e.\ each model $m \in \mathcal{M}$ induces a distribution $F^m$ over observables.

Let $F_0$ denote the observed data distribution. The identified set under assumption $\delta$ is $\mathcal{M}_0(\delta) = \{ m \in \mathcal{M}(\delta) : F^m = F_0 \}$


For a parameter of interest $\psi : \mathcal{M} \to \mathbb{R}$, the identified set is
\[
\Psi_0(\delta) = \{ \psi(m) : m \in \mathcal{M}_0(\delta) \}.
\]

As $\delta$ increases, the model class $\mathcal{M}(\delta)$ expands, so the identified set $\Psi_0(\delta)$ weakly enlarges. Thus $\delta$ can be interpreted as a \emph{sensitivity parameter} controlling the degree of identification.

\subsection{Connection to \textcite{nevo2012}}

We can see \textcite{nevo2012} imposing the following restrictions with $\delta = 1$:

\[
\operatorname{sign}\bigl(\operatorname{Cov}(Z,u)\bigr)
=
\operatorname{sign}\bigl(\operatorname{Cov}(X,u)\bigr),
\qquad
\lvert \operatorname{Corr}(Z,u) \rvert \le \delta \lvert \operatorname{Corr}(X,u) \rvert.
\]

These restrictions define a set of admissible data-generating processes consistent with observed moments and the structural parameter $\psi(m) = \beta$. This corresponds to specifying a model class $\mathcal{M}(\delta)$, where $\delta \in [0,1]$ .

\subsection{Identified Set in a Linear Model}
\subsubsection{Univariate instrument}

Redoing calculations from \textcite{nevo2012} we find the following bounds:

\[
\underline{\beta}(\delta)
=
\min\left\{
\beta^{IV}_{V(0)},
\beta^{IV}_{V(\delta)}
\right\},
\]
\[
\overline{\beta}(\delta)
=
\max\left\{
\beta^{IV}_{V(0)},
\beta^{IV}_{V(\delta)}
\right\}.
\]

See \autoref{app:proofs} for the full calculation.





\subsubsection{Bivariate instrument}

\section{General simulations}

\section{Networks}

\nocite{nevo2012}
\printbibliography


\appendix
\section{Identified Set in Linear Model}\label{app:proofs}

\[
Y=\beta X+U,\qquad E[U]=0.
\]

Assume
\[
\rho_{XU}\rho_{ZU}\ge 0,
\qquad
|\rho_{ZU}|\le \delta |\rho_{XU}|,
\qquad \delta\in[0,1].
\]

Define
\[
\lambda^*
=
\frac{\rho_{ZU}}{\rho_{XU}}.
\]

Then
\[
\lambda^*\in[0,\delta].
\]

For any \(\lambda\in[0,\delta]\), define
\[
V(\lambda)=\sigma_X Z-\lambda \sigma_Z X.
\]

The IV estimand using \(V(\lambda)\) is
\[
\beta^{IV}_{V(\lambda)}
=
\frac{\operatorname{Cov}(V(\lambda),Y)}
     {\operatorname{Cov}(V(\lambda),X)}.
\]

Equivalently,
\[
\beta^{IV}_{V(\lambda)}
=
\frac{
\sigma_X\sigma_{ZY}
-
\lambda\sigma_Z\sigma_{XY}
}{
\sigma_X\sigma_{XZ}
-
\lambda\sigma_Z\sigma_X^2
}.
\]

Hence the identified set is
\[
B(\delta)
=
\left\{
\beta^{IV}_{V(\lambda)}:\lambda\in[0,\delta]
\right\}.
\]

Since
\[
\frac{d}{d\lambda}
\beta^{IV}_{V(\lambda)}
=
\frac{
\sigma_X\sigma_Z
\left(
\sigma_{ZY}\sigma_X^2
-
\sigma_{XY}\sigma_{XZ}
\right)
}{
\left(
\sigma_X\sigma_{XZ}
-
\lambda\sigma_Z\sigma_X^2
\right)^2
},
\]
the sign of the derivative does not depend on \(\lambda\). Therefore the bounds are attained at the endpoints:
\[
\underline{\beta}(\delta)
=
\min\left\{
\beta^{IV}_{V(0)},
\beta^{IV}_{V(\delta)}
\right\},
\]
\[
\overline{\beta}(\delta)
=
\max\left\{
\beta^{IV}_{V(0)},
\beta^{IV}_{V(\delta)}
\right\}.
\]

Because
\[
V(0)=\sigma_X Z,
\]
we have
\[
\beta^{IV}_{V(0)}
=
\beta^{IV}
=
\frac{\sigma_{ZY}}{\sigma_{XZ}}.
\]

Thus
\[
\boxed{
B(\delta)
=
\left[
\min\left\{
\frac{\sigma_{ZY}}{\sigma_{XZ}},
\frac{
\sigma_X\sigma_{ZY}
-
\delta\sigma_Z\sigma_{XY}
}{
\sigma_X\sigma_{XZ}
-
\delta\sigma_Z\sigma_X^2
}
\right\},
\;
\max\left\{
\frac{\sigma_{ZY}}{\sigma_{XZ}},
\frac{
\sigma_X\sigma_{ZY}
-
\delta\sigma_Z\sigma_{XY}
}{
\sigma_X\sigma_{XZ}
-
\delta\sigma_Z\sigma_X^2
}
\right\}
\right].
}
\]

\end{document}
